\scanpage{155}
\mspage{100}
On reprend la description du ``foncteur'' $C$ (à valeurs dans $(\underline{\mathrm{Fib}}\,S)W_S^{-1}$, et alors sans guillemets), on trouve que $C(i_X)$, où
\[
X\xrightarrow{u=i_X}X'=C(X)
\]
est défini simplement (dans $(\underline{\mathrm{Fib}}_S)W_S^{-1}$) par ``la'' flèche $C(u)$ rendant commutatif le diagramme (dans $\underline{\mathrm{Fib}}\,S$)
\diagram{\begin{tikzcd}[ampersand replacement=\&]
X\arrow[r,"{u\,(=i_X)}"]\arrow[d,"i_X"'] \& \deleted{\unsure{C}}X'\,(=C(X))\arrow[d,"{i_{X'}\,(=i_{C(X)})}"]\\
C(X)\arrow[r,dash,"{C(u)(=C(i_X))}"] \& C(X')\,(=C(C(X))\text{,}
\end{tikzcd}}
de sorte qu'on a dans $\deleted{\unsure{F}}\mathcal M_{/S}$ le diagramme
\diagram{\begin{tikzcd}[ampersand replacement=\&]
X\arrow[r,shift left=.7ex,"{\substack{\mathrm{Cof}_0\cap W\\\unsure{\swarrow|\ni}\\i_X}}"]\arrow[r,shift right=.7ex,phantom,"{\deleted{\unsure{\longleftarrow|\longrightarrow}}}" description] \& \overset{\substack{\mathrm{Ob}\,\underline{\mathrm{Fib}}\,S\\\unsure{\swarrow|\ni}}}{C(X)}\arrow[r,shift left=.6ex,"i_{C(X)}"]\arrow[r,shift right=.6ex,"C(i_X)"'] \& \overset{\substack{\mathrm{Ob}\,\underline{\mathrm{Fib}}\,S\\\unsure{\downarrow|\ni}}}{C(C(X))}\text{,}
\end{tikzcd}}
commutatif i.e. $i_{C(X)}i_X=C(i_X)i_X$. Le Cor. page 92 nous dit que les flèches $i_{C(X)}$ et $C(i_X)$ sont égales dans $(\underline{\mathrm{Fib}}\,S)W_S^{-1}$.
Mais par hypothèse (p. 95), dans $\underline{\mathrm{Fib}}\,S$ on a $W_S=W_S^{\mathcal L}$, donc les flèches sont égales dans $(\underline{\mathrm{Fib}}\,S)(W_S^{\mathcal L})^{-1}$, et a fortiori dans $(\mathcal M_{/S})(W_S^{\mathcal L})^{-1}$, et on gagne.

Voyons ce qu'on peut faire dans le cas où on dispose d'un \uline{foncteur} $C$. On supposera bien sûr que ce foncteur dans $\mathcal M_{/S}$
% archive sheet number: 154
